\subsection{章末確認問題}

\begin{enumerate}[leftmargin=2em]
\item ソフトウェア技術者と単なるプログラマの役割の違いを説明せよ。
\item モジュール性、凝集性、疎結合性の違いを、身近な例で説明せよ。
\item コンピュータアーキテクチャとコンピュータ組織の違いを説明せよ。
\item ノイマン型とハーバード型の違いを説明せよ。
\item RISC と CISC の違いを、命令の複雑さと実行時間の観点から説明せよ。
\item 線形データ構造と非線形データ構造の違いを説明し、それぞれ例を二つ挙げよ。
\item 大O記法が実行時間そのものではなく増え方を表す理由を説明せよ。
\item サブルーチン、再帰、コルーチンの違いを説明せよ。
\item プロセスとスレッド、プロセス間通信、デッドロックの関係を説明せよ。
\item ACID と BASE の違いを、どのようなデータベースに向くかを含めて説明せよ。
\item OSI参照モデルの七層を順に挙げ、カプセル化とは何かを説明せよ。
\item AI for SE と SE for AI の違いを説明し、それぞれ例を二つ挙げよ。
\end{enumerate}

\begin{pointbox}{解答の方向性}
1は、実装だけでなく要求、設計、性能、テスト、保守まで扱う点を述べる。2は、分割、責務集中、依存低減を区別する。3は、外から見た構成と内部実現の違いを述べる。4から12は、本文の定義、表、注意欄の語を使って説明できる。

\end{pointbox}
